\subsection{EXAMPLES OF UNIVERSAL MAPPINGS}

* The examples which follow will, for the most part, be treated in detail elsewhere in this series.

I. Free algebraic structures. Let E be a set and let \(\Sigma\) be a species of algebraic structures, defined by one or more laws of composition. We take as morphisms the homomorphisms for the species \(\Sigma\) under consideration, and as \(\alpha\) -mappings arbitrary mappings of E into a \(\Sigma\) -set (in other words, \(\alpha\{x, s\} = \mathcal{F}(E, x)\) ). All the usual species of algebraic structures satisfy \((\mathrm{CU}_{\mathrm{III}})\) ; with the exception of division ring structures, they also satisfy \((\mathrm{CU}_{\mathrm{I}})\) , and \((\mathrm{CU}_{\mathrm{II}})\) is here a trivial consequence of \((\mathrm{CU}_{\mathrm{I}})\) .

Since in general there exist structures of species \(\Sigma\) defined on sets with at least two elements, the \(\alpha\) -mappings separate the elements of E, and E may therefore be considered as embedded in \(F_{E}\) . \(F_{E}\) is said to be the free \(\Sigma\) -set generated by E. Thus in algebra we speak of free semigroups, free groups, free modules, and free algebras.

\begin{enumerate}
\def\labelenumi{\Roman{enumi}.}
\setcounter{enumi}{1}
\item
  Rings and fields of fractions. Let E be a commutative ring with an identity element and let S be a multiplicatively closed subset of E which does not contain 0. We take \(\Sigma\) to be the species of structures of commutative rings with identity element, and the morphisms to be ring homomorphisms which transform identity element into identity element. The \(\alpha\) -mappings will be the homomorphisms \(\varphi\) of E into a commutative ring A with an identity element, such that \(\varphi(1)=1\) and \(\varphi(S)\) contains only units of A. The conditions \((\mathbf{Q}\mathbf{M}_{\mathbf{II}}), (\mathbf{C}\mathbf{U}_{\mathbf{I}})\) through \((\mathbf{C}\mathbf{U}_{\mathbf{III}})\) (with \(\mathfrak{a}=\operatorname{Card}(\mathbf{E})\operatorname{Card}(\mathbf{N}))\) are immediately verified. The universal mapping problem therefore always has a solution \((\mathbf{F}_{\mathbf{E}},\varphi_{\mathbf{E}})\) , but in general \(\varphi_{E}\) is not injective. The most frequently arising case is that in which E is an integral domain; in this case, \(\varphi_{E}\) is injective. If, moreover, we take \(S = E - \{0\}\) , then \(F_{E}\) is a field, called the field of fractions of E.
\item
  Tensor product of two modules. Let E be the product A × B of two modules over a commutative ring C which has an identity element. Take Σ to be the species of C-module structures, the morphisms to be linear mappings, and the α-mappings to be bilinear mappings of A × B into a C-module. The condition (QM \(_{II}\) ) is evidently satisfied, and so are (CU \(_{I}\) ) through (CU \(_{III}\) ) (with a = Card (E) Card (C) Card (N)). The universal C-module F \(_{E}\) corresponding to the pair (A, B) is called the tensor product of A and B and is written A ⊗ B. The universal mapping φ \(_{E}\) is written (x, y) → x ⊗ y; it is bilinear but not, in general, injective.
\item
  Extension of the ring of operators of a module. Let A be a commutative ring with an identity element, let B be a subring of A containing the identity element of A, and let E be a B-module. The species \(\Sigma\) is the species of A-module structures, the morphisms are A-linear mappings, and the \(\alpha\) -mappings are the B-linear mappings of E into an A-module. The universal A-module \(F_{E}\) corresponding to the B-module E is said to be obtained by extending to A the ring of operators B of E.
\end{enumerate}

V. Completion of a uniform space. Let E be a uniform space. Take \(\Sigma\) to be the species of structures of complete Hausdorff uniform spaces, the morphisms to be uniformly continuous mappings, and the \(\alpha\) -mappings to be uniformly continuous mappings of E into a complete Hausdorff uniform space. The \(\Sigma\) -admissible subsets of a complete Hausdorff uniform space are here the closed subsets (with respect to the topology of the space), and conditions \((\mathrm{QM}_{\mathrm{II}})\) and \((\mathrm{CU}_{\mathrm{I}})\) through \((\mathrm{CU}_{\mathrm{III}})\) are satisfied (with \(a = 2^{2^{\operatorname{Card}(E)}}\) ). The complete Hausdorff uniform space is (up to isomorphism) the completion of the Hausdorff uniform space associated with E (General Topology, Chapter II, §3, no. 7).

\begin{enumerate}
\def\labelenumi{\Roman{enumi}.}
\setcounter{enumi}{5}
\item
  Stone-Čech compactification. Let E be a completely regular space. \(\Sigma\) is the species of structures of compact spaces, the morphisms being continuous mappings (of a compact space into a compact space), and the \(\alpha\) -mappings are continuous mappings of E into a compact space. The \(\Sigma\) -admissible subsets are again the closed subsets, and conditions \((\mathrm{QM}_{\mathrm{II}})\) , \((\mathrm{CU}_{\mathrm{I}})\) through \((\mathrm{CU}_{\mathrm{III}})\) are easily verified (with the same cardinal as in Example V). The compact space \(F_{E}\) is (up to isomorphism) the ``Stone-Čech compactification'' obtained by completing E with respect to the coarsest uniformity for which all the continuous mappings of E into the interval {[}0, 1{]} of R are uniformly continuous (General Topology, Chapter IX, §1, Exercise 7); the mapping \(\varphi_{E}\) is injective, because any two distinct points of E can be separated by a continuous mapping of E into {[}0, 1{]}.
\item
  Free topological groups. Let E be a completely regular space, let \(\Sigma\) be the species of Hausdorff topological group structures, the morphisms being continuous homomorphisms; and take the \(\alpha\) -mappings to be the continuous mappings of E into a Hausdorff topological group. Conditions ( \(QM_{II}\) ) and ( \(CU_{I}\) ) to ( \(CU_{III}\) ) are easily verified, with
\end{enumerate}

\[
\mathfrak {a} = \text { Card } (\mathbf {E}) \text { Card } (\mathbf {N}).
\]

The Hausdorff topological group \(F_{E}\) which is the solution of this universal mapping problem for E is called the free topological group generated by the space E. Since any two distinct points of E can be separated by a continuous mapping of E into the Hausdorff topological group R, the mapping \(\varphi_{E}\) is injective; it can be shown that \(\varphi_{E}\) is a homeomorphism of E onto the subspace \(\varphi_{\mathrm{E}}(\mathrm{E})\) of \(F_{E}\) (*). Instead of taking \(\Sigma\) to be the species of structures of Hausdorff topological groups, we could also take other species of structures, such as those of Hausdorff topological abelian groups, compact groups, Hausdorff topological rings, Hausdorff topological vector spaces (over a topological division ring, considered as an auxiliary base set), etc.

\begin{enumerate}
\def\labelenumi{\Roman{enumi}.}
\setcounter{enumi}{7}
\item
  Almost periodic functions on a topological group. Let E be a topological group. Take \(\Sigma\) to be the species of compact group structures, the morphisms being continuous homomorphisms, and the \(\alpha\) -mappings being continuous homomorphisms of E into a compact group. Conditions (QM \(_{II}\) ), (CU \(_{I}\) ) through (CU \(_{III}\) ) are satisfied, with \(a = 2^{2^{\text{Card(E)}}}\) . The compact group F \(_{E}\) which is the solution of this universal mapping problem for E is called the compact group associated with E; the mapping \(\varphi_{E}\) is not necessarily injective. Every continuous real-valued function on E, of the form \(g \circ \varphi_{E}\) , where g is a continuous real-valued function on \(F_E\), is called an almost periodic function on E.
\item
  Albanese variety. Let E be an algebraic variety, let \(\Sigma\) be the species of structures of abelian varieties over the same base field as E (an abelian variety is a complete algebraic variety endowed with an algebraic group structure; it is necessarily commutative). The morphisms are rational mappings of one abelian variety into another (each morphism is necessarily the composition of a homomorphism and a translation). The \(\alpha\) -mappings are rational mappings of E into an abelian variety. Condition (CU \(_{I}\) ) is not satisfied, yet this universal mapping problem for E admits a solution \(F_{E}\) , called the Albanese variety of E. In general, the rational mapping \(\varphi_{E}\) is not injective. *
\end{enumerate}

